% sections/06_discussion.tex — Section VI (Discussion)
% Source: notes/paper_section_VI_draft.md v1 (Day 3 P1)
% Structure: §VI.1 alternative interpretations + §VI.2-5 four pushbacks
%            (Mantel setup, demo-supervised universality, auxiliary losses,
%             single benchmark) + §VI.6 future work + §VI.7 limitations.
% Converted Day 5 2026-05-13.

\section{Discussion}
\label{sec:discussion}

\subsection{Alternative Interpretations of the Findings}
\label{sec:discussion:alternatives}

The empirical findings of \S\ref{sec:results} admit alternative interpretations beyond the
one offered by Theorem~\ref{thm:bound}. We address the three most plausible.

\textbf{Alternative~1: Demonstration data is biased toward visually-distinctive tasks,
masking language-action correlation.} It is conceivable that LIBERO-Spatial demonstrations
exhibit near-zero language-trajectory correlation (M3) not because demonstrations are
intrinsically language-orthogonal, but because the 10 tasks happen to differ visually in
ways that dominate the inter-task distance matrix on trajectories. Under this
interpretation, a demonstration suite with visually-similar tasks differing only in
instruction semantics would show different M3 results.

This interpretation is partially testable. The mechanistic argument
(\S\ref{sec:intro:scope}, \S\ref{sec:discussion:single_benchmark}) predicts that
demonstration-trajectory correlation depends on the teleop protocol --- operators
executing motor commands based on visual perception, with language as label rather than
driver --- not on task visual distinctiveness. Within the present work, we cannot
conclusively distinguish this interpretation, since all our data comes from
LIBERO-Spatial. We note however that the backbone-capacity check
(\S\ref{sec:results:diagnostic}) rules out the related concern that high $\rho^2$ at the
backbone reflects visual rather than language structure: backbone-level $h_{\mathrm{OFT}}$
correlates with text-token-derived language embeddings (Gate~1 protocol), and the
F-statistic on action prediction confirms the backbone can produce task-discriminative
actions when explicitly probed.

\textbf{Alternative~2: The Gaussian MI surrogate underestimates the true MI, and the
actual demonstration-level language information is substantial.} The surrogate
$\hat{I} = -\tfrac{1}{2}\log(1 - \rho^2)$ assumes joint Gaussianity. If demonstration
trajectory and language are dependent through nonlinear or non-Gaussian structure, the
Mantel correlation $\rho$ underestimates the dependence, and $\hat{I}$ correspondingly
underestimates the true MI.

This is a real limitation acknowledged in \S\ref{sec:results:cfm_collapse}. We report the
squared correlation $\rho^2$ as the primary effect size, with the MI surrogate as
qualitative reference. The interpretation we offer rests on the \emph{gap} between
backbone-level and demonstration-level squared correlation ($8\times$ ratio), which is
invariant to monotone transformations of the underlying MI: if the Gaussian surrogate
underestimates true MI, it likely underestimates it at both backbone and demonstration
levels, preserving the ratio. The substantive claim --- that language conditioning
attenuates between backbone and demonstration --- is a relative observation that does not
depend on absolute MI estimates.

\textbf{Alternative~3: Theorem~\ref{thm:bound} is correct as stated but vacuous, because
no realistic policy achieves the bound.} It is logically possible that
demonstration-supervised policies all achieve $I(\Astar; \ell \mid o)$ far below the
bound, making the bound non-binding and uninformative.

Our empirical findings argue against this. The Gate~1 measurement shows the backbone has
$\rho_h^2 = 0.578$, equivalent to a Gaussian MI surrogate of $\sim 0.43$ nats. If the
bound were vacuous, policy-level language MI would substantially exceed the demonstration
bound of $\sim 0.04$ nats by leveraging the backbone's preserved language signal at no
cost. The C1 measurement on OTP-Soft V3 instead shows $r(z, \ell) = -0.19$ --- policy
latent inherits the demonstration-level near-orthogonality, not the backbone-level strong
correlation. The bound appears to be binding in practice.

\subsection{Pushback~1 --- Mantel Test Setup Validity}
\label{sec:discussion:pushback_mantel}

A skeptical reviewer may question the Mantel test methodology underpinning Gate~1, M3,
and C1: (a) the choice of distance metric (Euclidean) on potentially high-dimensional
representations; (b) the use of mean-pooled language embeddings as the language-side
distance matrix; (c) the 10-task LIBERO-Spatial set as a small-$N$ inter-task distance
test.

\textbf{(a) Distance metric choice.} Euclidean distance on per-task mean representations
is standard for representation-distance probing~\cite{kriegeskorte2008rsa}. The Mantel
test is invariant to monotone transformations of the distance matrices, so cosine
distance produces qualitatively identical patterns; we report Euclidean for consistency
with prior probing literature.

\textbf{(b) Language embedding choice.} We use mean-pooled token embeddings from the
OpenVLA-OFT tokenizer applied to the raw instruction prompt. This choice keeps the
language-side measurement aligned with what the backbone-internal text-processing
pipeline produces; alternative choices (CLIP text encoder, BERT pooling, instruction-tuned
sentence embeddings) would test slightly different language representations. The SC2
sanity check (\S\ref{sec:results:diagnostic}) verifies our pipeline is deterministic and
consistent between Gate~1 and C1 measurements.

\textbf{(c) Inter-task $N=10$.} The Mantel test with $N = 10$ tasks produces
$\binom{10}{2} = 45$ inter-task pairs and an effective sample size suitable for
permutation testing with $N_{\text{perm}} = 1000$. The replicated Gate~1 result ($\Delta
r = 0.0152$ between two independent runs of the same protocol on the same data) is well
within sampling noise for this $N$; we estimate sampling noise empirically rather than
relying on asymptotic distributional assumptions. The replication agreement supports the
robustness of the protocol at this sample size.

\subsection{Pushback~2 --- How Does Any VLA Architecture Achieve High SR?}
\label{sec:discussion:pushback_universality}

A natural objection to Theorem~\ref{thm:bound} is that it appears too strong: if all
demonstration-supervised architectures are bounded by demonstration-level language MI,
and demonstrations are near-language-orthogonal, then no demonstration-supervised policy
should perform well at language-conditioned tasks. Yet OFT achieves 97\% SR. How does
this fit?

The answer is in the two-scenario framing of \S\ref{sec:results:oft_contrast}. OFT's
97\% SR is achieved consistently with the bound: under scenario~(a), OFT carries small
but non-zero language MI within the bound, sufficient for task discrimination; under
scenario~(b), OFT carries zero language MI and achieves SR via observation-action mapping
alone. Both scenarios are consistent with the bound and the paraphrase invariance pattern
(\S\ref{sec:results:paraphrase}).

The conceptual error a reader might fall into is identifying ``task success'' with
``language conditioning use.'' Task success only requires the policy to produce correct
actions given the observation; it does not require the language instruction to causally
drive action selection. In LIBERO-Spatial (and many similar benchmarks), tasks are
visually distinctive --- each task has a unique object layout that can be discriminated
from observation alone. A policy that maps each visual scene to its correct action will
achieve high SR without using language. Theorem~\ref{thm:bound} bounds the
language-driven contribution to policy actions; it does not bound the SR-driven
contribution from observation processing.

We argue this explains why the standard SR + paraphrase-invariance evaluation regime
cannot detect the distinction. Both visually-grounded (b) and weakly-language-grounded
(a) policies produce the same SR profile under standard evaluation. Signal-flow
diagnostics (Gate~1, M3, C1) are needed to distinguish.

\subsection{Pushback~3 --- Auxiliary Language Losses as Bypass Route}
\label{sec:discussion:pushback_auxiliary}

We claim Theorem~\ref{thm:bound} is architecture-agnostic for any demonstration-supervised
policy, and that bypassing the bound requires interventions outside the demonstration
loss (Remark~\ref{rem:bypass}; \S\ref{sec:architecture:implications};
\S\ref{sec:architecture:falsification}). A reviewer may ask: what specifically counts as
``outside the demonstration loss,'' and can it actually achieve the bypass?

Several auxiliary objectives have been proposed in adjacent literature that, under our
framework, would modify the supervision distribution underlying
Theorem~\ref{thm:bound} and thereby alter the bound.

\textbf{Language reconstruction loss.} Training the policy to additionally reconstruct
the instruction from intermediate representations (style of~\cite{lynch2020multitask})
introduces a direct $\ell$-supervision signal that does not flow through demonstration
trajectories. Under our framework, this modifies the supervision distribution from
$\dataD(o, \ell, a)$ to a joint $\dataD(o, \ell, a) \cdot p_{\text{aux}}(\ell, z)$,
lifting the Theorem~\ref{thm:bound} bound on $I(\Astar; \ell \mid o)$ by the auxiliary
signal's information.

\textbf{Contrastive language-representation alignment.} Aligning policy hidden states
with backbone-derived language embeddings via contrastive objectives (CLIP-style
training, but for VLA hidden states) introduces a non-demonstration supervision pathway.
Backbone-preserved language information (Gate~1 $r = 0.76$) becomes available to the
policy through alignment rather than through demonstration mediation.

\textbf{Inference-time conditioning on backbone signals.} At test time, augmenting policy
inference with explicit backbone hidden state injection (analogous to retrieval-augmented
generation in NLP) bypasses the demonstration-trained policy's bound by drawing on the
backbone's preserved language signal directly.

These interventions are all \emph{outside} the demonstration-supervised loss in the sense
Theorem~\ref{thm:bound} targets. Our paper does not implement or empirically evaluate any
of them --- they are pointers to a research direction informed by
Theorem~\ref{thm:bound}, not solutions delivered by it.

A common potential confusion: does \emph{fine-tuning} the backbone on demonstrations
count as outside the demonstration loss? No --- fine-tuning incorporates the backbone
into the demonstration-supervised optimization, making it subject to
Theorem~\ref{thm:bound} just as the action head is. The Gate~1 measurement on the
\emph{frozen} OpenVLA-OFT backbone evidences the preserved language signal; fine-tuning
that backbone on demonstrations would degrade Gate~1 toward the demonstration-level
signal magnitude.

\subsection{Pushback~4 --- Single-Benchmark Scope}
\label{sec:discussion:single_benchmark}

A natural pushback to our diagnostic findings is single-benchmark scope: Gate~1
($r(h_{\mathrm{OFT}}, \ell) = 0.76$), M3 ($r \in [-0.27, -0.20]$), and the OFT contrast
(\S\ref{sec:results:oft_contrast}) are all demonstrated on LIBERO-Spatial. We separate
the theoretical and empirical components of our scope claim.

\textbf{Theoretical scope.} Theorem~\ref{thm:bound} bounds $I(\Astar; \ell \mid o) =
I_{\dataD}(a; \ell \mid o) \le \eps$. The right-hand side depends only on the joint
distribution of the supervision data $\dataD(\ell, a \mid o)$, not on the specific task
suite. The bound is benchmark-agnostic; the question is whether the antecedent
$I_{\dataD}(a; \ell \mid o) \approx 0$ holds across benchmarks.

\textbf{Empirical scope.} We expect demonstration-trajectory $\perp$ language to hold
broadly under standard teleoperated demonstration protocols, on a mechanistic argument:
teleoperators visually perceive the scene and execute motor commands, with language
serving as a post-hoc label affixed to the recorded trajectory rather than as a causal
driver of the recorded actions. Under this protocol, the joint distribution
$\dataD(\ell, a \mid o)$ factors approximately as $p(\ell \mid o) \cdot p(a \mid o)$,
yielding the language-orthogonality condition independently of task suite.

This mechanistic prediction is empirically testable. We provide one quantitative
instantiation (LIBERO-Spatial, $r = -0.27$); replication on LIBERO-Goal, LIBERO-Object,
and non-LIBERO teleoperated suites (e.g., RoboCasa, Bridge) is left to future work
(\S\ref{sec:discussion:future_work}). We do not claim universal cross-benchmark
generalization, but we do make a falsifiable mechanistic prediction: any teleop-collected
suite should yield $\lvert r(\text{traj}, \ell) \rvert \ll \lvert r(h_{\mathrm{backbone}},
\ell) \rvert$ under the same M3 protocol.

\subsection{Future Work}
\label{sec:discussion:future_work}

Three directions follow from this work.

\textbf{1.~Cross-benchmark empirical validation of M3.} Applying the M3 protocol to
LIBERO-Goal, LIBERO-Object, LIBERO-90, RoboCasa, and Bridge would test the mechanistic
prediction of \S\ref{sec:discussion:single_benchmark} across suites differing in task
structure, instruction complexity, and demonstration source. If the prediction holds, the
\S\ref{sec:discussion:single_benchmark} caveat strengthens to a general empirical claim
about teleop-collected VLA training data. If it fails on some suite, the failure would
itself be informative --- identifying which protocol features matter for
demonstration-level language conditioning.

\textbf{2.~Auxiliary-objective interventions to bypass Theorem~\ref{thm:bound}.} The
three intervention classes identified in \S\ref{sec:discussion:pushback_auxiliary}
(language reconstruction, contrastive alignment, inference-time conditioning) admit
concrete experimental designs. The most promising near-term direction is contrastive
alignment of policy hidden states with frozen backbone language embeddings --- this
leverages the Gate~1 result (backbone preserves language) without requiring new
demonstration collection. We expect this to materially lift the
Theorem~\ref{thm:bound} bound on policy-level language MI; the empirical question is
whether the resulting policy achieves higher language-driven SR distinguishable from
visually-grounded high SR under signal-flow diagnostics.

\textbf{3.~Diagnostic-methodology adoption in VLA evaluation.} Beyond the present paper's
findings, we propose that Gate~1, M3, and C1 measurements be reported alongside standard
SR + paraphrase-invariance results in future VLA evaluations. The three measurements are
cheap (each takes $<$~1 hour on a single GPU); they require only access to backbone
hidden states, demonstration trajectories, and policy latents, which are available for
nearly all published VLA architectures. Routine reporting would localize
language-conditioning behavior across the field and inform whether SR improvements
correspond to genuine language-conditioning gains or to better visual grounding.

\subsection{Limitations}
\label{sec:discussion:limitations}

We acknowledge five limitations explicitly:

\textbf{1.~Single benchmark.} \S\ref{sec:discussion:single_benchmark} discusses the scope
caveat in detail. The empirical findings rest on LIBERO-Spatial; cross-benchmark
validation is future work.

\textbf{2.~Single backbone.} Gate~1 measurements use the OpenVLA-OFT backbone
exclusively. Other VLAs (RT-2, $\pi_0$, generic LLM-based VLAs) might show different
Gate~1 magnitudes. We have not measured. The theoretical bound applies regardless of
backbone; only the empirical instantiation is backbone-specific.

\textbf{3.~Gaussian MI surrogate.} The quantitative MI estimates ($\hat{I}_h \approx
0.43$, $\hat{I}_{\mathrm{demo}} \approx 0.04$) rely on a Gaussian surrogate and should be
read as qualitative effect-size references rather than tight MI ceilings. The primary
effect-size finding is the squared-correlation ratio.

\textbf{4.~OTP-Soft architectural complexity.} OTP-Soft was originally designed under a
working hypothesis subsequently falsified by Theorem~\ref{thm:bound}
(\S\ref{sec:architecture:falsification}). It serves the present work as a diagnostic
platform rather than as an optimal architectural recommendation. Readers interested in
implementing the diagnostic methodology should adapt it to simpler VLA architectures
(OFT, RT-2, $\pi_0$) for which the architectural overhead is lower.

\textbf{5.~OFT scenario disambiguation.} \S\ref{sec:results:oft_contrast} leaves open
which of two scenarios (low-magnitude language preservation vs.\ entirely visual
grounding) describes OFT's high SR. Distinguishing requires either an OFT-specific
signal-flow diagnostic measurement (which we have not performed) or controlled
architecture experiments aligned with OFT's internal structure (future work).

These limitations bound the strength of our claims to: a theoretical bound with broad
applicability under stated conditions, one quantitative empirical instantiation, and a
diagnostic methodology proposed for broader adoption. The paper's contribution is
methodological and theoretical rather than a complete empirical resolution of the
language-conditioning question in VLAs.
